\chapter{Peripheral vascular disease}
\index{Peripheral vascular disease}

\medskip
\noindent\textit{Educational information; seek professional advice for individual care.}

\section*{What it usually means}
Peripheral vascular disease is a broad term for disease affecting blood vessels outside the heart and brain. It often refers to peripheral artery disease, in which narrowed arteries reduce blood flow to the legs, but vein and lymphatic conditions can cause different symptoms and treatments.

\section*{Possible symptoms}
Artery disease may cause exertional leg pain relieved by rest, cold or pale feet, numbness, weakness, slow-healing wounds, or pain at rest. Vein disease more often causes swelling, heaviness, skin changes, or ulcers. Sudden pain, swelling, colour change, or loss of function needs urgent assessment.

\section*{Assessment and management}
Assessment includes medical history, risk factors, pulses, skin and wound examination, and tests such as an ankle-brachial index or ultrasound. Treatment depends on the vessel involved and may include tobacco cessation, exercise, compression only when safe, wound care, medicines, or a procedure to restore flow. Do not start compression if significant arterial disease has not been assessed.

\section*{When to seek urgent care}
Seek emergency help for a suddenly cold, pale, blue, numb, weak, or severely painful limb; rapidly increasing swelling; chest pain or breathlessness with leg swelling; or a wound with spreading infection.

\section*{Sources}
\begin{itemize}
\item NHLBI peripheral artery disease diagnosis: \url{https://www.nhlbi.nih.gov/health/peripheral-artery-disease/diagnosis}
\item NHLBI peripheral artery disease treatment: \url{https://www.nhlbi.nih.gov/health/peripheral-artery-disease/treatment}
\end{itemize}
